\normalsize{Este anexo recoge los aspectos éticos, económicos, sociales y ambientales del
trabajo desarrollado, agrupados en los apartados siguientes.}
\vspace{0.1cm}

\Large{\textsc{\textcolor{nar}{B.1 Introducción}}} \par
\normalsize{Este trabajo se enmarca en el proyecto LINCE, una iniciativa del PERTE Aeroespacial para
el desarrollo de microsatélites LEO liderada por Indra. El objetivo es desarrollar el subsistema
de comunicaciones serie del ordenador de a bordo (OBC) del satélite, cubriendo tanto el
diseño de hardware electrónico (PCBs de prueba) como el firmware embebido sobre un
sistema operativo de tiempo real. Este tipo de tecnología tiene implicaciones relevantes en
ámbitos sociales, económicos, ambientales y éticos que se analizan a continuación.}
\vspace{0.3cm}

\Large{\textsc{\textcolor{nar}{B.2 Descripción de impactos relevantes relacionados con el proyecto}}} \par
\vspace{0.3cm}
\normalsize{%
Los impactos relevantes del trabajo se agrupan en cuatro ámbitos.

\textbf{Autonomía estratégica y tejido industrial (impacto económico y social).} LINCE se financia
a través del PERTE Aeroespacial con el objetivo declarado de que España disponga de capacidad
propia para producir en serie plataformas de microsatélite y sus cargas de pago. El subsistema de
comunicaciones del ordenador de a bordo es uno de los bloques que, de no desarrollarse aquí, habría
que adquirir a proveedores extracomunitarios sujetos a controles de exportación. Un transceptor
serie escrito en VHDL propio, con su código fuente disponible y su comportamiento medido y
documentado, puede auditarse, modificarse y volver a sintetizarse sin depender de licencias ni de
autorizaciones externas, a diferencia de un componente cerrado de terceros.

El desarrollo propio deja además conocimiento en el país. Un ingeniero que diseña y verifica un
transceptor aporta más valor por hora que quien integra un producto ya definido por otro proveedor,
y ese conocimiento (cómo se sintetiza un puente VHDL, cómo se caracteriza un driver serie a alta
velocidad) permanece aunque cambie el contratista. Comprar el mismo subsistema fuera no lo deja.

Formar ese personal cualificado lleva años y exige experiencia en programas reales, que solo se
adquiere participando en misiones. Retenerlo tampoco es sencillo: el perfil es internacionalmente
móvil y compite con agencias y fabricantes de países con industria espacial consolidada y con el
sector tecnológico civil. Sin continuidad de proyectos en España, el resultado habitual es la fuga
de titulados. El PERTE Aeroespacial no pretende sostener esa plantilla de forma indefinida con
fondos públicos, sino alcanzar una masa crítica de empresas y personal que permita a la industria
sostenerse después por contratos comerciales. Un trabajo de fin de máster dentro de un programa así
forma parte de ese mecanismo: aplica conocimiento académico a un problema real, con las
restricciones y los plazos de una misión concreta. Su resultado es alguien capacitado para seguir
trabajando en el sector.

\textbf{Doble uso civil y militar.} Una plataforma de microsatélites LEO con capacidad de
observación y de comunicaciones es tecnología de doble uso por definición: los mismos enlaces serie
que transportan telemetría de un sensor térmico sirven a una carga de pago de vigilancia. El
trabajo aquí desarrollado se sitúa en la capa más genérica posible del sistema (transporte de datos
entre periféricos y procesador) y no incorpora ninguna funcionalidad específica de aplicación, pero
sería deshonesto presentarlo como neutral: forma parte de un ordenador de a bordo cuyo uso final lo
determina la carga de pago que se le conecte y el operador que la explote. La responsabilidad sobre
ese uso recae en el marco regulatorio de exportación y en las decisiones del consorcio, no en el
componente, pero queda declarada.

\textbf{Residuos electrónicos y ciclo de vida del hardware.} El trabajo ha producido seis placas de
circuito impreso físicas (dos unidades de cada uno de los tres diseños), con más de doscientos
componentes montados y conectores metálicos, además de las revisiones descartadas durante el
desarrollo. Todo ello es residuo de aparatos eléctricos y electrónicos en el sentido de la Directiva
2012/19/UE~\cite{ue2012weee}. Los componentes seleccionados cumplen la Directiva 2011/65/UE de
restricción de sustancias peligrosas~\cite{ue2011rohs}, y por eso la lista de materiales especifica
encapsulados y acabados libres de plomo.

El proyecto adopta un enfoque NewSpace: componentes comerciales (COTS) en lugar de piezas
cualificadas para espacio. Esto reduce el coste y el plazo, pero acorta la vida útil esperada del
hardware y aumenta la frecuencia de sustitución, con el consiguiente aumento del volumen de residuo
por unidad de servicio prestado.

\textbf{Sostenibilidad del entorno orbital.} LINCE produce microsatélites de órbita baja destinados
a operar en constelación. La densidad de objetos en LEO es el principal problema de sostenibilidad
del sector, documentado anualmente por la Oficina de Basura Espacial de la ESA
\cite{esa_environment_report}, y toda decisión que acorte la vida operativa de un satélite (por
ejemplo, un fallo de aviónica no recuperable) adelanta su sustitución y, con ella, un lanzamiento
adicional y un objeto más en órbita.

\textbf{Consumo energético y presupuesto de potencia en órbita.} La lógica programable es
\textbf{el único elemento del MPSoC cuyo consumo depende directamente de las decisiones de diseño
tomadas en este trabajo}: la estimación de Vivado da 3,52~W para la variante GPIO, 3,86~W para la de
catorce controladores DMA y 3,62~W para la solución MCDMA finalmente seleccionada
(tabla~\ref{tab:bench_hw}). En un microsatélite de 100--200~kg, alimentado por paneles solares y
con un presupuesto de potencia cerrado, cada vatio de la aviónica es un vatio que no está
disponible para la carga de pago. Este impacto se analiza en detalle en el apartado siguiente.
}
\vspace{0.3cm}

\Large{\textsc{\textcolor{nar}{B.3 Análisis detallado de alguno de los impactos}}} \par
\vspace{0.3cm}
\normalsize{%
Se analiza en detalle el \textbf{impacto energético}, por ser el único de los anteriores sobre el
que este trabajo aporta medidas propias y sobre el que las decisiones de diseño tomadas tienen
efecto directo y cuantificable.

\textbf{Planteamiento.} La magnitud relevante en un satélite no es la potencia instantánea que
consume la aviónica, sino la \textbf{energía que cuesta entregar un byte útil}. Dos arquitecturas
que consumen lo mismo no son equivalentes si una transporta el triple de datos con esa potencia. La
comparativa del capítulo~4 permite calcular esa magnitud, porque mide sobre el mismo dispositivo,
el mismo transceptor y el mismo programa de pruebas tanto la potencia estimada por la herramienta
de implementación como el caudal realmente sostenido.

\textbf{Datos.} Con seis receptores simultáneos sobre el bus multipunto, la variante GPIO sostiene
20\,838~B/s con una potencia de 3,52~W y la variante MCDMA 68\,504~B/s con 3,62~W
(sección~\ref{sec:resultados_benchmark}). La energía por byte transportado resulta ser:

\begin{center}
\begin{tabular}{lrrr}
\hline
\textbf{Variante} & \textbf{Potencia} & \textbf{Caudal sostenido} & \textbf{Energía por byte} \\
\hline
GPIO  & 3,52~W & 20\,838~B/s & 169~$\mu$J \\
MCDMA & 3,62~W & 68\,504~B/s & 52,8~$\mu$J \\
\hline
\end{tabular}
\end{center}

La solución seleccionada consume 0,10~W más (un 2,8\,\%) y es \textbf{3,2 veces más eficiente por
byte transportado}. La variante intermedia de catorce controladores DMA, con 3,86~W, es además
la que más consume de las tres en términos absolutos.

\textbf{Interpretación.} El resultado no es una casualidad de la implementación: es la consecuencia
directa de haber quitado a la CPU de la ruta de datos. La arquitectura GPIO interrumpe al
procesador 1\,250 veces por kilobyte recibido, y cada una de esas interrupciones es un ciclo de
guardado y restauración de contexto que consume energía sin transportar información. El MCDMA baja
a 3 interrupciones por kilobyte. La lógica programable añadida (el puente VHDL y el controlador
multicanal) cuesta 0,10~W adicionales, pero libera a un procesador de cuatro núcleos que consume
bastante más que eso cuando está ocupado atendiendo interrupciones en lugar de en espera de
evento.

\textbf{Consecuencias sobre el sistema completo.} El margen liberado es doble y ambos efectos
tienen lectura ambiental. El primero es el margen de CPU: un procesador que no está atendiendo
interrupciones puede permanecer en espera de evento, lo que reduce el consumo medio del sistema de
procesamiento y, en un satélite alimentado por paneles solares, deja potencia disponible para la
carga de pago o permite dimensionar a la baja el generador fotovoltaico y las baterías, con la
consiguiente reducción de masa lanzada. El segundo es el margen de área: la solución seleccionada
ocupa el 5,18\,\% de la lógica del dispositivo frente al 14,67\,\% de la alternativa replicada, de
modo que su triplicación modular redundante para tolerancia a radiación mantiene el diseño en torno
al 16\,\% del dispositivo, mientras que la misma operación sobre la variante de catorce DMA
rondaría el 44\,\% y resultaría inviable. Ese margen es lo que hace posible endurecer el sistema
frente a los \textit{single event upsets} (SEU), y un satélite que sobrevive a un evento de radiación en
lugar de perderse es un satélite que no hay que sustituir: se evitan un lanzamiento y un objeto
adicional en una órbita ya congestionada~\cite{esa_environment_report}.

\textbf{Caudal de misión.} El cálculo anterior supone el canal trabajando a caudal sostenido. El
caudal real de la misión no está fijado todavía, y conviene acotar la conclusión sin adornar esa
falta de dato. Con tráfico alto por las interfaces, el argumento se sostiene por la energía por byte
y por el margen de CPU. Con tráfico bajo, si el canal no se satura, la ventaja en energía por byte
se diluye, pero la variante GPIO tampoco era una opción: ni por número de canales ni por coste de
interrupción. El cambio es favorable en los dos escenarios. Por eso la plataforma se dimensiona sin
fijar si las interfaces servirán a telemetría, a la carga de pago o a ambas, y el margen liberado,
de CPU y de área, vale mientras por ellas circule tráfico sostenido.

\textbf{Limitaciones del análisis.} Las cifras de potencia son la estimación del implementador de
Vivado, no una medida sobre el carril de alimentación de la placa, por lo que sirven para comparar
variantes entre sí, que es su propósito aquí, pero no como valor absoluto de consumo. El cálculo
tampoco incluye el consumo del sistema de procesamiento, que es donde el efecto descrito debería ser
mayor.
}
\vspace{0.3cm}

\Large{\textsc{\textcolor{nar}{B.4 Conclusiones}}} \par
\vspace{0.3cm}
\normalsize{%
Desde el punto de vista \textbf{económico}, el trabajo contribuye a un objetivo declarado de
política industrial: sustituir por desarrollo propio un bloque del ordenador de a bordo que de otro
modo se adquiriría fuera de la Unión Europea. El coste directo del desarrollo, detallado en el
Anexo~\ref{anx:presupuesto}, es de un orden de magnitud inferior al de la adquisición y cualificación de una solución
comercial equivalente, y el resultado queda en el laboratorio como base reutilizable.

Desde el punto de vista \textbf{social}, el impacto es indirecto pero real: la plataforma forma
parte de una constelación de observación y comunicaciones cuya explotación afecta a la capacidad de
vigilancia y a los servicios de comunicación disponibles. El trabajo se sitúa en una capa
deliberadamente genérica del sistema, pero se declara aquí su condición de tecnología de doble uso
en lugar de presentarlo como un componente neutral.

Desde el punto de vista \textbf{medioambiental}, el balance tiene dos signos. En el lado negativo,
el trabajo genera residuo electrónico (seis placas fabricadas más las revisiones descartadas) y se
inscribe en un enfoque NewSpace basado en componentes comerciales que acorta la vida útil del
hardware y aumenta la frecuencia de sustitución. En el lado positivo, la arquitectura de transporte
seleccionada es 3,2 veces más eficiente por byte transportado que la de partida y deja margen de
área suficiente para endurecer el diseño frente a radiación, lo que alarga la vida operativa
esperada del satélite y reduce la necesidad de lanzamientos de reposición. La selección de
componentes se ha ceñido a la Directiva RoHS~\cite{ue2011rohs}, y las placas descartadas deben
gestionarse como residuo de aparatos eléctricos y electrónicos a través del canal habilitado por el
laboratorio, conforme a la Directiva RAEE~\cite{ue2012weee}.

Desde el punto de vista \textbf{ético y de gestión}, el criterio seguido a lo largo de la memoria
ha sido declarar las limitaciones en lugar de omitirlas: las tarjetas se identifican explícitamente
como hardware de validación funcional y no de vuelo, se documenta la interfaz que quedó sin validar
(SpaceWire) y el motivo, se declaran las limitaciones de cada método de medida y se señalan los
errores de diseño detectados y sus correcciones. Las decisiones de arquitectura se sostienen sobre
medidas propias y reproducibles, con el código y los bancos de prueba publicados
\cite{repo_tfm}, y no sobre estimaciones. Ese es el compromiso con la responsabilidad profesional
que el ejercicio de la ingeniería exige y que este anexo pretende acreditar.
}
